\clearpage
\section{Introduction}

\subsection{Motivation}
In the modern embedded systems development landscape, the cost of software exceeds the cost of designing the hardware it runs on~\cite{niespodziany2025}. As semiconductor technology advances and customer demand for broader, more sophisticated functionality grows, an embedded application is assembled by dozens of engineers across several teams, each one producing independent modules that recombine into many product variants across the device's lifetime.

The industry's response to this assembly problem is to formalize it: separate the modules from the application that consumes them, and describe their composition in a structure independent of any single module's source code. The \textit{diff} framework, presented by Niespodziany~\cite{niespodziany2025}, embodies that response for C++ embedded software through data-driven dependency injection: the application's composition is described in a JSON topology file consumed at startup.

While \textit{diff} improves reusability and flexibility, composing a product variant means editing the topology file manually. The developer must ensure its structural correctness, with limited tool support. This flow does not meet the demands of a reliable, repeatable engineering process~\cite{humble2010continuous}, and these demands grow with the size of the system.

The industry's assembly problem and the framework's own proposal~\cite{niespodziany2025} together motivate \textbf{\textit{diff-forge}}~: a visual tool that supports managing the dependency graph, and producing JSON file consumed by the existing C++ without modification. Next section sets out the specific gaps that \textit{diff-forge} aims to resolve.

\subsection{Problem Statement}
Manual maintenance of a \textit{diff} topology file imposes three major gaps along the developer's workflow. Each grows with the size of the system and the number of developers who work with the file.

\textbf{Lack of module discovery.}
Modules are distributed across Artifactory repositories, with many modules per repository and many versions per module. To compose a new system, the developer must learn which modules exist, what interfaces they expose, and what parameters they accept. This gap scales with the size of the catalog and the number of new systems assembled.

\textbf{Lack of graph readability.}
The dependency graph is encoded as text. To trace which components depend on a given module, or whether one component depends on another directly or transitively, the developer reads the JSON or runs a custom query against it. Every architectural question, from the impact of a change to the presence of a cycle, requires this reconstruction.

\textbf{Lack of real-time validation feedback.}
The JSON parser only checks syntax and is not capable of detecting lexical errors within each entry's configuration fields (such as identifiers, dependency references, parameters, \ldots), or structural errors across the JSON file (cycles, unmet dependencies, interface mismatches). These errors are caught only when the binary starts, in CI, or during code review. A broken build blocks other developers, and responsibility for fixing it often falls to someone other than the original author.

\subsection{Project Goals and Scope}
The primary goal of this project is to develop \texttt{diff-forge}, a cross-platform desktop application that enables developers to visually compose, verify, and export software topologies. The application utilizes modern web technologies including TypeScript, React~\cite{reactjs}, and Electron~\cite{electronjs}. The specific goals of the application include:
\begin{itemize}
    \item Providing an intuitive drag-and-drop interface for instantiating components from a managed catalog.
    \item Implementing a real-time validation engine to prevent structural errors before they reach the production environment.
    \item Ensuring seamless integration with existing C++ development workflows by exporting framework-compliant configuration files.
\end{itemize}

\subsection{Thesis Structure}
The following chapters comprise this thesis: Chapter~2 reviews the state of the art in dependency injection and visual programming. Chapter~3 provides an overview of the \texttt{diff} framework and its target use cases. Chapter~4 describes the architectural design of \texttt{diff-forge}, while Chapter~5 details its implementation. Lastly, Chapter~6 evaluates the tool's effectiveness, followed by concluding remarks in Chapter~7.
